\section{The \texttt{Pol\_guide\_vmirror} McStas Component}
Polarising guide with nvs x 2 supermirrors sitting in v-shapes inside,
upgraded from original single V-cavity.

\subsection*{Identification}
\begin{itemize}
  \item \textbf{Site:} 
  \item \textbf{Author:} Peter Willendrup
  \item \textbf{Origin:} DTU
  \item \textbf{Date:} June 2022
\end{itemize}

\subsection*{Description}
\begin{lstlisting}
Models a rectangular guide with entrance centered on the Z axis and
with nvs x two supermirrors sitting in v-shapes inside.
The entrance lies in the X-Y plane.  Draws a true depiction
of the guide with mirrors, and trajectories.
The polarisation is handled similar to in Monochromator_pol.
The reflec functions are handled similar to Pol_mirror.
The up direction is hardcoded to be along the y-axis (0, 1, 0)

The reflectivity parameters can be
1. double pointer initializations with 5 parameters (e.g. {R0, Qc, alpha, m, W} AND inputType=0), or
2. double pointer initializations with 6 parameters (e.g. {R0, Qc, alpha, m, W, beta} AND inputType=0), or
3. table names (e.g."supermirror_m2.rfl" AND inputType=1), or
4. individually assigned values (e.g. rR0=1.0, rQc=0.0217, ... etc AND inputType=2).
NB! This might cause warnings by the compiler that can be ignored.
Note that the reflectivity functions that use with values and those that use tables are different.

GRAVITY: YES

%BUGS
No absorption by mirror.
\end{lstlisting}

\subsection*{Input parameters}
Parameters in \textbf{boldface} are required; the others are optional.

\begin{longtable}{p{0.22\textwidth}p{0.12\textwidth}p{0.46\textwidth}p{0.14\textwidth}}
\toprule
\textbf{Name} & \textbf{Unit} & \textbf{Description} & \textbf{Default} \\
\midrule
\endhead
nvs & 1 & Number of V-cavities across width of guide & 1 \\
xwidth & m & Width at the guide entry & 0.1 \\
yheight & m & Height at the guide entry & 0.1 \\
length & m & length of guide & 0.5 \\
rR0 & 1 & Cavity wall reflectivity below critical scattering vector, use with inputType = 2 & 1.0 \\
rQc & AA-1 & Cavity wall reflectivity critical scattering vector, use natural Ni (m=1 definition) value of 0.0217, use with inputType = 2 & 0.0217 \\
ralpha & AA & Cavity wall slope of reflectivity, use with inputType = 2 & 6.5 \\
rmSM & 1 & Cavity wall m-value of material. Zero means completely absorbing. Use with inputType = 2 & 2 \\
rW & AA-1 & Cavity wall width of reflectivity cut-off, use with inputType = 2 & 0.00157 \\
rbeta & AA2 & Cavity wall curvature of reflectivity, use with inputType = 2 & 80 \\
rUpR0 & 1 & Mirror spin up reflectivity below critical scattering vector, use with inputType = 2 & 1.0 \\
rUpQc & AA-1 & Mirror spin up reflectivity critical scattering vector, use natural Ni (m=1 definition) value of 0.0217, use with inputType = 2 & 0.0217 \\
rUpalpha & AA & Mirror spin up slope of reflectivity, use with inputType = 2 & 2.47 \\
rUpmSM & 1 & Mirror spin up m-value of material. Zero means completely absorbing. Use with inputType = 2 & 4 \\
rUpW & AA-1 & Mirror spin up width of reflectivity cut-off, use with inputType = 2 & 0.0014 \\
rUpbeta & AA2 & Mirror spin up curvature of reflectivity, use with inputType = 2 & 0 \\
rDownR0 & 1 & Mirror spin down reflectivity below critical scattering vector, use with inputType = 2 & 1.0 \\
rDownQc & AA-1 & Mirror spin down reflectivity critical scattering vector, use natural Ni (m=1 definition) value of 0.0217, use with inputType = 2 & 0.0217 \\
rDownalpha & AA & Mirror spin down slope of reflectivity, use with inputType = 2 & 1 \\
rDownmSM & 1 & Mirror spin down m-value of material. Zero means completely absorbing. Use with inputType = 2 & 0.65 \\
rDownW & AA-1 & Mirror spin down width of reflectivity cut-off, use with inputType = 2 & 0.003 \\
rDownbeta & AA2 & Mirror spin down curvature of reflectivity, use with inputType = 2 & 0 \\
debug & 1 & if debug \textgreater{} 0 print out some internal runtime parameters & 0 \\
allow\_inside\_start & 1 & Allow neutron to start inside cavity (no propagation to entry plane) & 0 \\
rPar & 1 & Cavity wall parameters array for rFunc, use with inputType = 0 & \{1.0, 0.0217, 6.5, 2, 0.00157, 80\} \\
rUpPar & 1 & Mirror Parameters array for rUpFunc, use with inputType = 0 & \{1.0, 0.0217, 2.47, 4, 0.0014\} \\
rDownPar & 1 & Mirror Parameters for rDownFunc, use with inputType = 0 & \{1.0, 0.0217, 1, 0.65, 0.003\} \\
rParFile &  & Cavity wall parameters filename for rFunc, use with inputType = 1 & "" \\
rUpParFile &  & Mirror Parameters filename for rUpFunc, use with inputType = 1 & "" \\
rDownParFile &  & Mirror Parameters filename for rDownFunc, use with inputType = 1 & "" \\
\bottomrule
\end{longtable}

\subsection*{Links}
\begin{itemize}
  \item \href{run:/home/willend/willend-McCode/mcstas-comps/optics/Pol_guide_vmirror.comp}{Source code} for \texttt{Pol\_guide\_vmirror.comp}.
\end{itemize}
\IfFileExists{Pol_guide_vmirror_static.tex}{\input{Pol_guide_vmirror_static.tex}}{}